\chapter*{KATA PENGANTAR}
\addcontentsline{toc}{chapter}{KATA PENGANTAR}

Puji syukur penulis panjatkan ke hadirat Tuhan Yang Maha Esa karena atas rahmat dan karunia-Nya, penulis dapat menyelesaikan Laporan Tugas Akhir yang berjudul "Perancangan dan Implementasi Warehouse Management Controller \& Sistem Lokalisasi untuk Prototipe Smart Warehouse". Laporan ini disusun sebagai salah satu syarat untuk menyelesaikan program Sarjana di Program Studi Sistem dan Teknologi Informasi, Sekolah Teknik Elektro dan Informatika, Institut Teknologi Bandung.

Dalam proses penyelesaian tugas akhir ini, penulis mendapatkan banyak bantuan dan dukungan dari berbagai pihak. Oleh karena itu, penulis ingin menyampaikan ucapan terima kasih kepada:

\begin{enumerate}
    \item Bapak Dr.techn. Ary Setijadi Prihatmanto, S.T., M.T. selaku dosen pembimbing 1 tugas akhir yang membimbing kelompok tugas akhir penulis pada topik Smart Warehouse. Penulis mengucapkan terima kasih atas segala arahan, bimbingan, dan saran yang diberikan dalam proses perancangan dan penulisan laporan ini.

    \item Bapak I Gusti Bagus Baskara Nugraha, S.T., M.T., Ph.D. selaku dosen pembimbing 2 tugas akhir yang membimbing penulis dalam proses penyusunan tugas akhir yang telah memberikan masukan dan saran dalam penyempurnaan Tugas Akhir ini.

    \item Timotius Vivaldi Gunawan, Muhammad Shafa Sauditya, dan Rubi Naufal Tiandito, rekan satu kelompok tugas akhir penulis. Perjalanan bersama, sejak penyusunan proposal topik hingga rampungnya produk tugas akhir ini, tidak akan berjalan selancar ini tanpa kerja sama dan kekompakan kalian.

    \item Ibu dan Nenek penulis, yang doa serta dukungannya tidak pernah putus. Terima kasih telah berjuang tanpa lelah untuk mempersiapkan bekal kehidupan penulis, baik hari ini maupun di masa yang akan datang.

    \item Teman-teman K2 (Athar, Nicho, Favian, Raihan) yang telah mewarnai tiga tahun masa perkuliahan penulis di Program Studi Sistem dan Teknologi Informasi. Kebersamaan kalian menjadi salah satu bagian paling berkesan dari perjalanan ini.

    \item Mas Fahmi dan Mas Vitra, asisten Pak Ary, atas kesediaannya membimbing dan membantu penulis selama pengerjaan tugas akhir ini.
\end{enumerate}

Penulis menyadari bahwa laporan ini masih jauh dari sempurna. Oleh karena itu, kritik dan saran yang membangun sangat diharapkan untuk perbaikan di masa mendatang. Semoga laporan Tugas Akhir ini dapat memberikan manfaat bagi pembaca dan pihak-pihak terkait.

\begin{flushright}
    Bandung, \DTMlangsetup{showdayofmonth=false,showmonthname=true,showyear=true}\today\\

    Penulis \\
    \vspace{1cm}
    Raizan Iqbal Resi
\end{flushright}
